% P42 | 5.6 Independent validators | Result: least-squares Monte Carlo and the quasi-analytic approximation
% guide: dev_notes/paragraph_guides/P42_result_least_squares_monte_carlo_and_the_quasi_a.txt
Least-squares Monte Carlo works backward on simulated paths: at each exercise date it regresses the discounted future cash flow on a polynomial of the spot, using only in-the-money paths, and exercises when the payoff exceeds the fitted continuation value.
The estimate is biased low, because the regression policy is not the optimal one, and it approaches the Bermudan value as the basis grows.
With $50$ exercise dates and $2\cdot10^5$ paths the estimate is \LSMfifty\ with standard error \LSMfiftySE, against \BermFifty\ for a Bermudan tree on the same dates.
The difference is within one standard error, as it is for the other date counts in Table~\ref{tbl:lsm}.
The last column of Table~\ref{tbl:lsm} shows that the gap between the American and the Bermudan value halves as the number of exercise dates doubles.
Restricting exercise to fifty dates therefore costs about $0.0006$ in value, which is below the standard error of the regression estimate.
The Barone--Adesi--Whaley approximation, which approximates the early-exercise premium by a power of the spot, gives \BAWV\ at the baseline, an error of \BAWerr, and a critical price of \BAWSc\ against the tree boundary of about \AmBearly\ at the start.
Table~\ref{tbl:baw} shows that its absolute error grows with the maturity and reaches about one cent for the three-year put, which is why it serves for quick estimates and not for final prices.
